\begin{definition} $2\sat = \{\varphi \ | \text{ $\varphi$~--- выполнимая в 2-КНФ формула}\}$. \end{definition}

\begin{theorem} $2\sat \in \nlcsp$. \end{theorem}

\textbf{TODO} Нужно предъявить сводимость к PATH

\begin{proof}
\begin{enumerate}
    \item Покажем принадлежность $\nlsp$, для этого решим $2\sat$ через $\PATH$. Пусть $\varphi$ в 2-КНФ и зависит от $x_1, \ldots, x_n$. Воспользуемся некоторым метод из матлога о взаимнообратном отображении булевой формулы в орграф. А именно:
    
    Построим граф $G_\varphi$ с $2n$ вершинами, которые будут обозначаться $x_1, \ldots, x_n, \overline{x}_1, \ldots, \overline{x}_n$. Рассмотрим, какие будут рёбра при каких дизъюнктах (для простоты переменные: x, y):
    \begin{enumerate}
        \item Если $\varphi$ содержит дизъюнкт $(x \lor y)$, то добавим в граф рёбра $\overline{x} \to y$ и $\overline{y} \to x$.
        \item Если $\varphi$ содержит дизъюнкт $(\overline{x} \lor y)$, то добавим в граф рёбра $x \to y$ и $\overline{y} \to \overline{x}$.
        \item Если $\varphi$ содержит дизъюнкт $(\overline{x} \lor \overline{y})$, то добавим в граф рёбра $x \to\overline{y}$ и $y\to\overline{x}$.
    \end{enumerate}
    
    Тогда одновременное отсутствие путей из $x$ в $\overline{x}$ и из $\overline{x}$ в $x$ ($x$ и $\overline{x}$ не в одной КСС) равносильно выполнимости формулы:
    
    \begin{itemize}
        \item $\Rightarrow$ Раскрасим граф в два цвета, где если $x$ имеет истинный цвет, то $\overline{x}$ окрашена в ложный. Тогда для выполняющего набора верно, что если начало ребра окрашено в истину, то и его конец тоже, что невозможно, если есть пути из $x$ в $\overline{x}$ и из $\overline{x}$ в $x$.
        
        \item $\Leftarrow$ Пусть нет пути из $x$ в $\overline{x}$, тогда не может быть одновременно путей $x \to y$ и $x \to \overline{y}$ (иначе есть путь $\overline{y} \to \overline{x}$, а тогда есть путь $x \to y \to \overline{x}$, что неверно). Сделаем $x$ и все достижимые из него литералы истинными, это можно сделать непротиворечиво.
    \end{itemize}
    
    Теперь, если в дизъюнкте $(z \lor t)$ $z$ стал ложью, то есть путь $x \to \overline{z}$, а еще есть ребро $\overline{z} \to t$ по построению, тогда $t$ достижим, и он истинен. Значит в каждой скобке либо хотя бы одна ложь, либо оба неизвестны (это легко решается тем, что эти переменные могут быть любыми).

    Получается, свели $2\sat$ к проверке отсутствия нескольких путей, а это задача из $\nlsp$. Тогда в сертификате проверяем все $2n$ путей в порядке возрастания.

    \item Теперь покажем полноту, для этого сведем $\overline{\PATH}$ к $2\sat$. Построим по тройке $(G, s, t)$ формулу так: для каждой вершины заводим переменную, а для ребра $u \to v$ дизъюнкт $(\overline{u} \lor v)$ и заведём дизъюнкты для $s$ и для $t$. Если пути из $s$ в $t$ нет, то рассмотрим набор, в котором все достижимые из $s$ переменные истинны, а остальные~--- ложны. Формула выполнима, так как все дизъюнкты при таком наборе будут истинными ($\overline{t}$ истинен, так как $t$ не достижима, а в остальных не может быть комбинации 10, так как если достижима $u$, то достижима и $v$).  А если формула выполнима, то любая переменная, достижимая из $s$ истинна (индукция по длине пути), но $\overline{t}$ тоже истинно, а значит $t$ недостижима из $s$.
\end{enumerate}

\end{proof}
